\subsubsection{最初の小小節}

ここに本文を書く。
主テキストの順序を保ちつつ，必要な定義，文脈，証明の行間を自然な本文として補う。

\begin{dfn}[定義 1.1]\notelabel{dfn:1.1}
ここに定義を書く。
\end{dfn}

\begin{prop}[命題 1.2]\notelabel{prop:1.2}
ここに命題を書く。
\end{prop}

\begin{proof}
ここに証明を書く。
証明が長い場合は，主テキストと手書きノートから証明段階を読み取り，自然な段落と接続文で分ける。

まず次の主張を示す。

\begin{claim}\notelabel{claim:1.2.local}
ここに証明内の局所的な主張を書く。
\end{claim}

この主張は，定義に戻って必要な条件を確認すれば従う。
ここに局所的な主張の証明を書く。\localqed

以上から命題が従う。
\end{proof}

\begin{rmk}
既に示した内容を再利用する場合は，同じ主張を再証明せず，たとえば \dfnref{1.1} や \claimref{1.2.local} のように参照する。
\end{rmk}
